\documentclass[11pt,a4paper]{article}
\usepackage[utf8]{inputenc}
\usepackage[spanish,es-noshorthands]{babel}
\usepackage{amsmath,amssymb}
\usepackage{geometry}
\geometry{top=2.5cm, bottom=2.5cm, left=2.5cm, right=2.5cm}
\usepackage{xcolor}
\usepackage{tcolorbox}
\tcbuselibrary{skins, breakable}
\usepackage{hyperref}
\usepackage{enumitem}
\usepackage{listings}
\usepackage{siunitx}
\usepackage{circuitikz}
\usepackage{microtype}

% Definición de colores
\definecolor{azulOscuro}{RGB}{20, 45, 85}
\definecolor{azulClaro}{RGB}{235, 243, 250}
\definecolor{grisBorde}{RGB}{180, 200, 220}

\hypersetup{
    colorlinks=true,
    linkcolor=azulOscuro,
    urlcolor=azulOscuro,
    pdftitle={Monitor de Salud IoT con ESP32 y Dashboard Web}
}

\lstset{
    basicstyle=\ttfamily\small,
    breaklines=true,
    frame=single,
    rulecolor=\color{grisBorde},
    backgroundcolor=\color{azulClaro},
    literate=
      {á}{{\'a}}1 {é}{{\'e}}1 {í}{{\'i}}1 {ó}{{\'o}}1 {ú}{{\'u}}1
      {Á}{{\'A}}1 {É}{{\'E}}1 {Í}{{\'I}}1 {Ó}{{\'O}}1 {Ú}{{\'U}}1
      {ñ}{{\~n}}1 {Ñ}{{\~N}}1 {₂}{{$_2$}}1 {µ}{{$\mu$}}1
}

\title{\textbf{\Large Monitor de Salud IoT con ESP32 y Dashboard Web}\\[4pt]
\large Especificación y Prompt de Desarrollo Paso a Paso}
\author{\textbf{Proyecto Electrónica \& IoT}}
\date{\today}

\begin{document}

\maketitle

\begin{tcolorbox}[colback=azulClaro, colframe=azulOscuro, title=\textbf{Descripción del Modelo / Prompt}, arc=2mm, fonttitle=\bfseries]
Este documento contiene la especificación detallada en forma de prompt estructurado para guiar el desarrollo completo de un \textbf{Monitor de Salud del Paciente basado en ESP32 con Dashboard Web Local en tiempo real}, abarcando el ensamblaje físico, conexiones de sensores, firmware C++ e interfaz web.
\end{tcolorbox}

\vspace{0.3cm}

\section*{Prompt Sugerido para LLM}

\begin{description}[style=nextline, leftmargin=1cm]
    \item[\textbf{Rol:}] Actúa como un experto en sistemas embebidos, electrónica médica e IoT.
    
    \item[\textbf{Objetivo:}] Generar una guía de implementación paso a paso para construir un \textbf{Monitor de Salud del Paciente basado en ESP32 con Dashboard Web Local en tiempo real}.
    
    \item[\textbf{Restricción importante:}] Omite cualquier paso relacionado con el diseño o fabricación de placas de circuito impreso (PCB). Todo el ensamblaje físico debe explicarse utilizando una protoboard, módulos comerciales y conexiones con cables jumper.
\end{description}

\hrulefill

\vspace{0.3cm}

\section*{Estructura de la Guía de Implementación}

\begin{enumerate}[label=\textbf{\arabic*.}, leftmargin=*]

    \item \textbf{Lista de Componentes y Función:}
    \begin{itemize}
        \item \textbf{Microcontrolador:} ESP32 Dev Module.
        \item \textbf{Módulos y Sensores:}
        \begin{itemize}
            \item Módulo ECG AD8232 con sus electrodos.
            \item Oxímetro de pulso MAX30102 para medición de $\text{SpO}_2$ y ritmo cardíaco.
            \item Sensor ambiental BME280 para temperatura y humedad ambiente.
            \item Módulo ADC de 16 bits ADS1115 (comunicación $\text{I}^2\text{C}$).
            \item Termistor NTC de $10\,\text{k}\Omega$ para temperatura corporal.
            \item Capacitores de $10\,\mu\text{F}$ y $100\,\text{nF}$ para filtrado y estabilización de la alimentación.
        \end{itemize}
    \end{itemize}

    \item \textbf{Esquema de Conexiones en Protoboard:}
    \begin{itemize}
        \item Explicar la conexión del bus $\text{I}^2\text{C}$ en los pines \textbf{GPIO 21 (SDA)} y \textbf{GPIO 22 (SCL)} del ESP32, compartido por el ADS1115, BME280 y MAX30102.
        \item Conexión de la salida analógica del AD8232 al canal \textbf{AIN0} del ADS1115 y sus pines de control digital (\texttt{SDN}, \texttt{LO+}, \texttt{LO-}) al ESP32.
        \item Conexión del divisor de voltaje del termistor NTC a la entrada \textbf{AIN1} del ADS1115.
        \item Colocación de los capacitores de desacoplo cerca de las líneas de alimentación para reducir ruido en la señal de ECG.
    \end{itemize}

    \item \textbf{Firmware C++ (Arduino IDE):}
    \begin{itemize}
        \item Librerías necesarias a instalar (\texttt{ADS1115}, \texttt{BME280}, \texttt{MAX30105}, \texttt{WiFi}) y selección de la placa ESP32 Dev Module.
        \item Lógica del código: muestreo continuo de la señal ECG a $\approx 125\,\text{Hz}$, procesamiento de buffers de $\text{SpO}_2$ y lecturas cada 2 segundos de temperatura y humedad.
        \item Configuración del servidor web local y creación del endpoint \texttt{/api/vitals} para responder en formato JSON.
        \item Inclusión del archivo de cabecera \texttt{dashboard\_page.h} con el código HTML, CSS y JavaScript de la interfaz web.
    \end{itemize}

    \item \textbf{Colocación de Sensores en el Paciente:}
    \begin{itemize}
        \item Posicionamiento correcto de los 3 electrodos del AD8232 en el cuerpo.
        \item Fijación adecuada del sensor MAX30102 en el dedo y del termistor NTC sobre la piel.
    \end{itemize}

    \item \textbf{Prueba y Uso del Dashboard:}
    \begin{itemize}
        \item Apertura del Monitor Serie a 115200 baudios para obtener la dirección IP asignada por la red Wi-Fi.
        \item Explicación de las secciones de la interfaz web:
        \begin{itemize}
            \item \textbf{Monitor de Signos Vitales:} Gráfico ECG en tiempo real, BPM, $\text{SpO}_2$ y temperaturas.
            \item \textbf{Registro de Paciente}.
            \item \textbf{Exportación de datos a Excel}.
            \item \textbf{Alertas Clínicas}.
        \end{itemize}
    \end{itemize}

\end{enumerate}

\vspace{0.3cm}

\section*{Estimación de Costos del Hardware}

Para la construcción del prototipo, se ha recopilado la estimación aproximada de precios de mercado minorista (en USD) para los componentes requeridos:

\begin{center}
\small
\begin{tabular}{|l|l|c|}
\hline
\textbf{Componente} & \textbf{Función / Descripción} & \textbf{Costo Aprox. (USD)} \\ \hline
ESP32 Dev Module & Microcontrolador Wi-Fi / Bluetooth & \$6.00 -- \$10.00 \\ \hline
AD8232 ECG Module & Sensor ECG con electrodos y cableado & \$6.00 -- \$12.00 \\ \hline
MAX30102 & Oxímetro de pulso ($\text{SpO}_2$ y ritmo cardíaco) & \$6.00 -- \$10.00 \\ \hline
BME280 & Sensor ambiental (temperatura y humedad) & \$4.00 -- \$10.00 \\ \hline
ADS1115 & ADC de 16-bits $\text{I}^2\text{C}$ (4 canales analógicos) & \$4.00 -- \$6.00 \\ \hline
Termistor NTC 10K & Sensor analógico de temperatura corporal & \$0.50 -- \$1.50 \\ \hline
Accesorios Prototipado & Protoboard, cables jumper y capacitores desacoplo & \$5.00 -- \$10.00 \\ \hline
\hline
\textbf{Costo Total Estimado} & \textbf{Prototipo completo en protoboard} & \textbf{\$31.50 -- \$59.50 USD} \\ \hline
\end{tabular}
\end{center}

\textit{Nota: Los costos indicados son valores de referencia basados en distribuidores de componentes electrónicos comerciales. Pueden variar según el fabricante, impuestos locales y gastos de envío.}

\vspace{0.3cm}

\section*{Requerimiento de Salida}
Proporciona la explicación estructurada, clara y con ejemplos de código en C++ y estructuras JSON donde sea necesario.

\end{document}
